% !TeX encoding = UTF-8
\documentclass[variant=apuntes,cover=institutional,mode=screen]{ua-engineering-v6-article}
% !TeX encoding = UTF-8
\UASetUniversity{Universidad Austral}
\UASetFaculty{Facultad de Ingeniería}
\UASetDepartment{Departamento de Computación}
\UASetCampus{Pilar}
\UASetCareer{Ingeniería Informática}
\UASetCommission{Comisión A}
\UASetProfessor{Equipo docente}
\UASetSemester{Segundo cuatrimestre 2026}
\UASetCourse{Sistemas Distribuidos}
\UASetDate{2026}


\UASetClassNumber{04}
\UASetClassTitle{Consenso distribuido}
\UASetDocKind{Apuntes}

\title{Clase 04 --- Consenso distribuido: fundamentos, imposibilidad y algoritmos}
\author{\UAProfessor}
\date{\UADate}

\begin{document}
\maketitle

\section{Introducción}

En la clase anterior estudiamos replicación y consistencia. Vimos que múltiples réplicas introducen divergencia debido a asincronía, fallas y concurrencia.

Esto nos lleva a un problema central:

\begin{quote}
¿Cómo imponemos un orden único de operaciones en un sistema distribuido?
\end{quote}

La respuesta es el \textbf{consenso distribuido}.

\begin{uakeybox}
Consenso permite que múltiples nodos acuerden un único valor, lo que induce un orden global de operaciones.
\end{uakeybox}

---

\section{Modelo del sistema}

Trabajamos bajo el siguiente modelo:

\begin{itemize}
\item $n$ procesos (nodos)
\item comunicación por mensajes
\item red asíncrona (sin límites de tiempo)
\item fallas por crash (fail-stop)
\end{itemize}

No asumimos:

\begin{itemize}
\item relojes sincronizados
\item tiempos máximos de entrega
\end{itemize}

---

\section{Definición formal de consenso}

Cada proceso $p_i$ propone un valor $v_i$.

El algoritmo debe garantizar:

\subsection{Agreement}
\begin{quote}
Si dos procesos deciden, deciden el mismo valor.
\end{quote}

\subsection{Validity}
\begin{quote}
El valor decidido fue propuesto por algún proceso.
\end{quote}

\subsection{Termination}
\begin{quote}
Todo proceso correcto eventualmente decide.
\end{quote}

---

\section{Imposibilidad: Teorema FLP}

\subsection{Enunciado}

\begin{quote}
En un sistema completamente asíncrono con al menos una posible falla, no existe algoritmo determinista que garantice consenso (terminación).
\end{quote}

---

\subsection{Intuición profunda}

El problema no es Agreement ni Validity, sino \textbf{Termination}.

La red puede:

\begin{itemize}
\item retrasar mensajes arbitrariamente
\item simular fallas (delay vs crash indistinguibles)
\end{itemize}

---

\subsection{Estados de ejecución}

Definimos:

\begin{itemize}
\item \textbf{Univalente}: el resultado está determinado (0 o 1)
\item \textbf{Bivalente}: el resultado aún puede ser 0 o 1
\end{itemize}

El estado inicial es típicamente bivalente.

---

\subsection{Idea de la prueba}

1. Existe un estado inicial bivalente  
2. Siempre existe una transición que mantiene bivalencia  
3. El scheduler puede elegir esas transiciones indefinidamente  

\begin{quote}
Resultado: ejecución infinita sin decisión
\end{quote}

---

\subsection{Interpretación}

\begin{uakeybox}
No se puede garantizar consenso perfecto sin introducir supuestos adicionales (timeouts, sincronía parcial, etc.)
\end{uakeybox}

---

\section{Consenso en la práctica}

Los sistemas reales evaden FLP mediante:

\begin{itemize}
\item timeouts
\item elecciones probabilísticas
\item sincronía parcial
\end{itemize}

Esto permite progreso \emph{en la práctica}, aunque no en el peor caso teórico.

---

\section{Raft: consenso práctico}

Raft es un algoritmo diseñado para ser comprensible y seguro.

---

\subsection{Componentes}

\begin{itemize}
\item Leader
\item Followers
\item Log replicado
\end{itemize}

---

\subsection{Estructura temporal}

El tiempo se divide en \textbf{términos} (terms).

\begin{itemize}
\item cada término tiene un líder
\item los términos son monótonos
\end{itemize}

---

\subsection{Elección de líder}

\begin{enumerate}
\item follower expira timeout
\item se convierte en candidate
\item solicita votos
\item si obtiene mayoría $\to$ leader
\end{enumerate}

---

\subsection{Replicación de log}

\begin{enumerate}
\item cliente envía request al leader
\item leader agrega entrada a su log
\item leader replica a followers
\item cuando mayoría confirma $\to$ commit
\end{enumerate}

---

\section{Invariantes de Raft}

Estos garantizan \textbf{seguridad}.

---

\subsection{Log Matching Property}

\begin{quote}
Si dos logs tienen una entrada en el mismo índice y término, entonces todos los elementos anteriores son idénticos.
\end{quote}

\textbf{Razón:}
\begin{itemize}
\item el leader solo replica desde prefijos consistentes
\item followers rechazan inconsistencias
\end{itemize}

---

\subsection{Leader Completeness}

\begin{quote}
Si una entrada está comprometida, aparecerá en todos los líderes futuros.
\end{quote}

\textbf{Idea:}
\begin{itemize}
\item commit requiere mayoría
\item futuros líderes también requieren mayoría
\item intersección $\Rightarrow$ preservación
\end{itemize}

---

\section{Quorums}

Un quorum es un subconjunto de nodos que puede tomar decisiones.

Para consenso:

\[
|Q| > \frac{N}{2}
\]

---

\subsection{Propiedad clave}

\[
Q_1 \cap Q_2 \neq \emptyset
\]

\textbf{Demostración:}

Supongamos dos quorums disjuntos:

\[
|Q_1| + |Q_2| \le N
\]

Pero:

\[
|Q_1| > N/2, \quad |Q_2| > N/2
\]

\[
|Q_1| + |Q_2| > N
\]

Contradicción.

---

\section{Paxos}

Paxos es el algoritmo clásico de consenso.

---

\subsection{Roles}

\begin{itemize}
\item Proposer
\item Acceptor
\item Learner
\end{itemize}

---

\subsection{Fases}

\textbf{Phase 1 (Prepare)}
\begin{itemize}
\item proposer envía número $n$
\item acceptors prometen no aceptar menor
\end{itemize}

\textbf{Phase 2 (Accept)}
\begin{itemize}
\item proposer propone valor
\item acceptors aceptan si cumple promesa
\end{itemize}

---

\section{Propiedad de seguridad de Paxos}

\begin{quote}
Si un valor es elegido, ningún otro valor distinto puede ser elegido.
\end{quote}

---

\subsection{Prueba (idea)}

\begin{itemize}
\item elección requiere mayoría
\item nuevas propuestas deben consultar mayoría
\item intersección $\Rightarrow$ conocimiento previo
\end{itemize}

Por lo tanto:

\begin{quote}
no puede aparecer un valor distinto
\end{quote}

---

\section{Consenso y linearizabilidad}

El consenso permite implementar:

\begin{itemize}
\item orden total de operaciones
\item máquina de estados replicada
\end{itemize}

\begin{uakeybox}
Consenso $\Rightarrow$ orden total $\Rightarrow$ linearizabilidad
\end{uakeybox}

---

\section{Costo del consenso}

\begin{itemize}
\item múltiples rondas de comunicación
\item necesidad de mayoría
\item latencia
\item menor disponibilidad bajo partición
\end{itemize}

---

\section{Modelo del curso}

\[
D = (P, F, G, S, T)
\]

\begin{itemize}
\item $P$: estado consistente global
\item $F$: fallas y asincronía
\item $G$: consenso (agreement)
\item $S$: Raft/Paxos
\item $T$: latencia y disponibilidad
\end{itemize}

---

\section{Conclusión}

\begin{quote}
Consenso es el mecanismo que transforma un sistema distribuido en una máquina secuencial lógica.
\end{quote}

Pero:

\begin{itemize}
\item no es gratis
\item no siempre progresa
\item requiere compromisos explícitos
\end{itemize}

\begin{uakeybox}
El consenso no elimina la incertidumbre: la controla.
\end{uakeybox}


\section{Síntesis razonada de la clase}
Esta unidad debe leerse como parte de una cadena de decisiones de diseño y no como un catálogo de mecanismos. Para estudiar el tema con profundidad conviene reconstruir cuatro preguntas: \emph{qué problema aparece}, \emph{qué garantía se desea}, \emph{qué mecanismo la aproxima} y \emph{qué costo introduce}. En cada ejemplo debe identificarse además el modelo de falla implícito.

\begin{uakeybox}
Una respuesta técnicamente madura no termina al nombrar un algoritmo o patrón: declara sus supuestos, la propiedad que preserva y el escenario en el que deja de ser suficiente.
\end{uakeybox}

\section{Bibliografía guiada y ruta de profundización}
Para profundizar esta clase se recomienda: Ongaro y Ousterhout (Raft); Lamport (Paxos); Fischer, Lynch y Paterson (FLP); Kleppmann, cap. 9.

La lectura debe hacerse con preguntas concretas: (1) ¿qué modelo de sistema asume el autor?, (2) ¿qué propiedad demuestra o persigue?, (3) ¿qué falla o anomalía motiva el mecanismo?, y (4) ¿qué trade-off queda abierto?

\section{Conexión con el Trabajo Final Integrador}
El grupo debe señalar qué parte de su sistema utiliza los conceptos de esta unidad, justificar por qué son necesarios y documentar una alternativa descartada. Cuando el contenido todavía no corresponda a la etapa vigente, debe registrarse como hipótesis o decisión pendiente.

\section{Autoevaluación conceptual}
\begin{enumerate}
\item ¿Cuál es el problema distribuido concreto que motiva esta clase?
\item ¿Qué propiedad de seguridad y qué propiedad de progreso están en juego?
\item ¿Qué supuesto de red, tiempo o falla resulta decisivo?
\item ¿Qué trade-off aceptaría en un sistema real y cuál no aceptaría en un dominio crítico?
\item ¿Cómo observaría experimentalmente que la solución se comporta como espera?
\end{enumerate}

\section{Profundización autocontenida v9}
\subsection{Problema, límite e invariante}
El núcleo de esta clase es consenso, FLP, mayorías, Raft, Paxos, seguridad, progreso y máquina de estados replicada. Para estudiarlo de manera autónoma conviene separar tres planos. Primero, el \emph{problema observable}: qué comportamiento incorrecto puede ver un cliente o un operador. Segundo, el \emph{límite del modelo}: qué información, sincronía o coordinación no está disponible. Tercero, el \emph{invariante}: qué propiedad no debe violarse aunque el sistema pierda progreso temporalmente.

El modelo de fallas relevante incluye crash, demora arbitraria, elección dividida y pérdida temporal de mayoría. Una solución debe describirse como protocolo y no solo como nombre de tecnología: actores, estados, mensajes, condición de éxito, condición de aborto y estado visible después de una interrupción. La evidencia esperada es términos, votos, logs, commitIndex y un invariante de seguridad.

\subsection{Método de análisis}
Ante un caso nuevo, siga esta secuencia:
\begin{enumerate}
\item Identifique actores, dato o recurso compartido y frontera de confianza.
\item Escriba una ejecución nominal y una ejecución adversarial.
\item Distinga seguridad (lo malo no ocurre) de progreso (algo bueno finalmente ocurre).
\item Declare qué supuesto permite avanzar: mayoría, deadline, sincronía parcial, operación idempotente, almacenamiento durable u otro.
\item Compare una alternativa más simple y una más fuerte; registre qué métrica o experimento haría cambiar la decisión.
\end{enumerate}

\subsection{Caso transferible}
Considere una plataforma de reservas que recibe operaciones duplicadas, pierde conectividad entre regiones y debe sostener un camino crítico sin ocultar el estado real al usuario. Aplique el mecanismo de esta clase a una única decisión concreta. Una respuesta madura no intenta resolver todo: delimita la operación, formula la garantía y explica la degradación esperada cuando el supuesto central deja de cumplirse.

\subsection{Errores de razonamiento frecuentes}
\begin{itemize}
\item confundir una propiedad con el producto que suele implementarla;
\item describir el camino feliz sin recorrer una falla intermedia;
\item afirmar ``exactly-once'', ``alta disponibilidad'' o ``consistencia'' sin definir alcance observable;
\item medir solo throughput aunque la hipótesis trate sobre semántica o recuperación;
\item presentar la complejidad añadida como beneficio sin compararla con la alternativa más simple.
\end{itemize}

\section{Lectura guiada}
Ongaro y Ousterhout, Raft; Lamport, Paxos Made Simple; material MIT 6.5840 sobre consenso.

Durante la lectura, registre: modelo de sistema, propiedad perseguida, contraejemplo que motiva el mecanismo, costo principal y una condición bajo la cual no lo elegiría. Estas cinco anotaciones convierten la bibliografía en una herramienta de diseño y no en una lista para memorizar.

\section{Aplicación al Trabajo Final Integrador}
El equipo debe producir un ADR breve con la decisión asociada a esta clase. Debe contener: contexto, dos alternativas, supuesto decisivo, garantía elegida, trade-off, evidencia pendiente y fecha de revisión. La evidencia mínima esperada para esta unidad es: términos, votos, logs, commitIndex y un invariante de seguridad.

\section{Autoevaluación avanzada}
\begin{enumerate}
\item ¿Puedo explicar la clase sin nombrar primero una tecnología?
\item ¿Puedo construir una ejecución que viole la garantía si el mecanismo se configura mal?
\item ¿Puedo distinguir seguridad, progreso y experiencia del usuario?
\item ¿Puedo proponer una medición que valide la propiedad, no solo el rendimiento?
\item ¿Puedo defender cuándo una solución más simple sería suficiente?
\end{enumerate}

\end{document}
